\newpage
\section*{Willkommen zur FOSSGIS-Konferenz 2023 in Berlin!}\label{welcome}
Die Abkürzung { FOSSGIS} steht für {\bfseries f}reie und {\bfseries O}pen"={\bfseries S}ource"={\bfseries S}oftware für {\bfseries G}eo{\bfseries i}nformations{\bfseries s}ysteme.
Die FOSSGIS-Konferenz 2023 wird vom gemeinnützigen FOSSGIS e.V., der
OpenStreetMap"=Community und der Humboldt-Universität zu Berlin
veranstaltet.
Ziel der jährlich stattfindenden Konferenz ist die Verbreitung von freier,
quelloffener Software für Geoinformationssysteme. In den nächsten vier Tagen
haben Sie die Gelegenheit, sich mit Entwickler:innen und anderen Anwender:innen
auszutauschen und \mbox{neueste} Informationen zu Anwendungen und
Arbeitsmöglichkeiten zu erhalten.

\section*{Anwender- und Expert:innentreffen, Birds of a Feather}
Die FOSSGIS-Konferenz ist eine Communityveranstaltung.
Die gesamte Konferenz über stehen zwei Räume für Anwender- und Expert:innentreffen oder spontan organisierte
Treffen Gleichgesinnter (Birds of a Feather), u.\,ä.
zur Verfügung. Eigene Sessions können Sie an der Pinnwand im Foyer selbst eintragen.
\pagebreak

\section*{Abendveranstaltung am Mittwoch}\label{schwaetzli}
Traditionell findet am ersten Abend der FOSSGIS-Konferenz die
Abendveranstaltung statt, auch Social Event genannt. Der Eintritt
ist im FOSSGIS-Konferenz-Ticket enthalten. Eine Anmeldung ist erforderlich.
Die Abendveranstaltung wird am Konferenzstandort im Erwin-Schrödinger-Zentrum
ab {\bfseries 18:30 bis 22:00 Uhr} stattfinden.

\section*{Mapathon mit Ärzte ohne Grenzen am Freitag und Samstag}
Ärzte ohne Grenzen lädt im Rahmen der FOSSGIS Konferenz 2023 zu einem Mapathon ein. Die Initiative hat das Ziel die Erstellung und Vervollständigung von Kartenmaterial zu unterstützen, welches die Hilfe für Menschen in Krisengebieten erleichtern kann. Es wird eine Einführung geben. Weitere Informationen sind hier zu finden. FOSSGIS-Teilnehmende sind eingeladen, mitzumachen, auch wenn sie noch nie gemappt haben.
Termine: {\bfseries Freitag, 17.00 bis 20.00 Uhr} und {\bfseries Samstag, 11.00 bis 14.00 Uhr}. Bitte melden Sie sich über den Anmeldelink auf der Konferenzwebseite an.
\pagebreak

\section*{Exkursionen}
\subsection*{Kartenabteilung der Staatsbibliothek zu Berlin}
Von den Sektionskarten von Balbi und der Preußischen Uraufnahme zu digitalen Kartendaten.
Am {\bfseries Freitagnachmittag} und {\bfseries Samstagvormittag} werden Führungen durch die Karten\-abteilung der Staatsbibliothek zu Berlin im Haus Unter den Linden angeboten. Anmeldung erforderlich.

\subsection*{Campus Adlershof/WISTA-Gelände}
Die FOSSGIS 2023 findet auf dem Gelände des Wissenschafts- und Wirtschaftsstandort Adlershof (WISTA) statt. Dort existieren noch mehrere technische Denkmale, wie beispielsweise der Trudelturm, die an den Beginn der Luftfahrtindustrie auf dem Campus erinnern. Zu sehen sind aber auch moderne Forschungseinrichtungen, Gründer- und Technologiezentren, die auf die heutige Entwicklung als Zukunftsort für Forschung und Innovation verweisen. In einer kleinen, etwa 60-minütigen Führung können wir sowohl dieser Geschichte als auch den neuen Trends nachstöbern und mehr über den Standort Adlershof erfahren.
Die Exkursion startet am {\bfseries Freitag} um {\bfseries 16:30 Uhr} (nach dem Sektempfang) vor dem Haupteingang des Erwin-Schrödinger-Zentrums. Anmeldung erforderlich.

\section*{Rahmenprogramm am Donnerstag}
\subsection*{Gruppenfoto}
Auch in diesem Jahr wollen wir uns das Gruppenfoto nicht entgehen lassen und laden Sie ein am {\bfseries Donnerstag} in der {\bfseries Nachmittagspause} zum gemeinsamen Gruppenfoto am Haupteingang des Erwin-Schrödinger-Zentrums ein.

\subsection*{Treffen der Freiberufler}
Am {\bfseries Donnerstag} werden sich um {\bfseries 17:45 Uhr} die Freiberufler aus dem FOSSGIS-Bereich im Seminarraum 1\textquotesingle 306 zum gemeinsamen Austausch treffen.

\subsection*{Mitgliederversammlung des FOSSGIS e.\,V.}
Am {\bfseries Donnerstag} sind alle Mitglieder und Gäste ab {\bfseries 18.00 Uhr} herzlich eingeladen, an der Mitgliederversammlung teilzunehmen und sich zu beteiligen. Der FOSSGIS e.\,V. lädt zum Diskutieren, Kennenlernen, Abstimmen und zu Neuwahlen ein. Es wird Getränke und Pizza für alle geben. Der Verein freut sich über zahlreiches Erscheinen.

\section*{Rahmenprogramm am Freitag}
\subsection*{Jeopardy-Quiz}
Am {\bfseries Freitag} gibt es vor dem Abschluss des reguläres Konferenzprogramms um {\bfseries 14.30 Uhr} das nunmehr legendäre und sehr humorvolle Quiz in Form des FOSSGIS-Jeopardy mit Hannes und Tobia.

\subsection*{Sektempfang am FOSSGIS-Stand}
Alle Mitglieder des FOSSGIS-Vereins, Freunde und Interessierte sind am {\bfseries Freitag} ab {\bfseries 16.00 Uhr} herzlich zum Sektempfang zum Ausklang der FOSSGIS 2023 am FOSSGIS-Vereins-Stand eingeladen.

\subsection*{OSM-Event am Freitagabend}
Für alle, die am {\bfseries Freitagabend} noch in der Stadt sind und/oder am OSM-Event teilnehmen möchten gibt es um {\bfseries 19/19:30 Uhr} ein gemeinsames Treffen im Restaurant Greek Food Olympia, Rudower Chaussee 5a in 12489 Berlin-Adlershof. Jeder zahlt seine Rechnung selbst.

\small
\newpage
\label{platinsposoren}
\cleardoubleevenpage
\section*{Platinsponsor und Aussteller}
\begin{center}
  \centerline{\includegraphics[width=1.0\textwidth]{001_WhereGroup.jpg}}
\end{center}
\vspace*{-0.4cm}

\noindent
Die {\bfseries WhereGroup GmbH} ist eines der größten Open-Source-Soft\-ware\-häuser der GEO-IT-Branche in Deutschland. Wir managen zahlreiche komplexe GIS-Projekte unterschiedlichster Art. Als mittelständisches Unternehmen mit über 40 Mitarbeiter*innen an vier Standorten arbeiten wir innovativ und haben dennoch die Bodenhaftung nicht verloren.

\noindent
Wir bieten Ihnen kompetente Unterstützung in den Bereichen Geographische Informationssysteme (GIS), Web-GIS, Datenbanken, Standards, Interoperabilität und System-Integration. Angefangen bei der Beratung, Konzeption und Entwicklung bis hin zum Betrieb dynamischer Kartenanwendungen im Intra- und Internet.
\newpage
\noindent
Grundlage unseres Schaffens ist der Open-Source-Gedanke. Als Teil einer starken Community, mit der wir in engem Austausch stehen, engagieren wir uns aus Überzeugung im FOSSGIS e.\,V., im QGIS-DE e.\,V. und bei der OSGeo. Wir sind aus voller Überzeugung auf Open-Source-Entwicklungen spezialisiert und integrieren professionelle freie Software nahtlos in proprietäre Systeme.

\noindent
Wir implementieren Lösungen für den Mittelstand, die Industrie und auf allen Ebenen der öffentlichen Verwaltung, beraten und begleiten bei Planung, Migration und Einführung von raumbezogenen Informationssystemen.

\noindent
Dabei reicht das Spektrum unserer Projekte von Desktop\-Lösungen über Geoportalen und kartenbasierter Datenverwaltung bis hin zu hochverfügbaren Anwendungen für die freie Wirtschaft und die öffentliche Verwaltung.

\noindent
Unser Schulungsinstitut, die FOSS Academy, bietet außerdem praxisorientierte Schulungen zum Thema "`GIS mit Open-Source-Software"' an.

\noindent
Mehr zur WhereGroup unter www.wheregroup.com und www.foss-academy.com.


\newpage
\section*{Platinsponsor und Aussteller}
%\begin{center}
  \includegraphics[width=1.0\textwidth]{002_camptocamp.png}
  %\end{center}
  \vspace{1.0\baselineskip}
  
\noindent
{\bfseries CamptoCamp} Camptocamp gehört zu den führenden Dienstleistern im Bereich Open-Source-GIS und ist in vielen unterschiedlichen Open-Source-Communitys stark engagiert.
\newpage
\noindent
Unsere Dienstleistungen stützen sich auf 20 Jahre Erfahrung in der Umsetzung von innovativen GIS-Lösungen für Behörden und Unternehmen und erlauben einen hochwertigen und individuellen Service. Das Besondere an Camptocamp sind die hochqualifizierten Mitarbeiter und ihr großes Engagement im "Ökosystem" der eingesetzten Open-Source Software-Lösungen, indem sehr enge Beziehungen zu den Herstellern der jeweiligen Produkte gepflegt werden.

\noindent
Um die oft anspruchsvollen Projekte umzusetzen, erstellt Camptocamp individuelle Lösungen, die auf den am besten geeigneten und fortschrittlichsten Open Source-Technologien basieren. Camptcamp ist in München, Lausanne, Olten, Paris und Chambéry vertreten und bietet neben Lösungen im GIS-Bereich auch eine große Expertise im den Bereichen ERP (Enterprise-Resource-Planning) und IT-Infrastruktur-Lösungen.

\newpage
\cleardoubleevenpage
\section*{Platinsponsor und Aussteller}
%\vspace*{-0.7\baselineskip}
%\begin{center}
  \includegraphics[width=0.7\textwidth]{003_dataport.png}
  %\end{center}
  \vspace{1.0\baselineskip}
  
\noindent
    {\bfseries Dataport - Anstalt des öffentlichen Rechts}
    \vspace{1.0\baselineskip}
    
\noindent
    {\bfseries Das Unternehmen}
    
\noindent    
Dataport ist der Partner für die Digitalisierung des öffentlichen Sektors. Für Schleswig-Holstein, Hamburg, Bremen und Sachsen-Anhalt. Und für die Steuerverwaltung in Niedersachsen und Meck\-lenburg-Vorpommern. Als IT-Dienstleister gestaltet Dataport den digitalen Wandel gemeinsam mit Ländern und Kommunen. Mit rund 5.000 Mitarbeiter:innen an acht Standorten. Das Unternehmen erzielte 2022 einen Umsatz von 1,18 Milliarden Euro.    
Dataport ist eine Anstalt des öffentlichen Rechts. Träger sind die Länder Schleswig-Holstein, Hamburg, Bremen, Sachsen-Anhalt, Niedersachsen und Meck\-lenburg-Vorpommern. Das Unternehmen arbeitet nicht gewinnorientiert, sondern strebt in enger Absprache mit den Trägern ein ausgeglichenes Betriebsergebnis an
\vspace{1.0\baselineskip}

\noindent
{\bfseries Der Auftrag}

\noindent
Dataport stellt dem öffentlichen Sektor alle benötigten IT-Services zur Verfügung. Dazu gehören der Betrieb von Infrastrukturen wie Rechenzentrum, Netze und Clients oder die zentrale Beschaffung von Informationstechnologie (IT). Außerdem die Entwicklung und der Betrieb von Software. Dataport unterstützt bei allen Aspekten der Digitalisierung. Durch umfassendes Consulting, Projektmanagement, Innovationsmanagement oder Geschäftsprozessmanagement.
\newpage
\noindent
{\bfseries Die Werte}
\vspace{1.0\baselineskip}

\noindent
{\em Digitale Sicherheit:}\\
Das Twin Data Center ist eines der sichersten Rechenzentren in Europa. Zertifiziert vom Bundesamt für Sicherheit in der Informationstechnik (BSI) und vom TÜViT. Der Betrieb von IT-Systemen folgt eindeutigen Prozessen und Regeln unter ständiger Kontrolle durch ein professionelles Sicherheitsmanagement. Das sorgt für Sicherheit.
\vspace{1.0\baselineskip}

\noindent
{\em Digitale Souveränität:}\\
Damit der Staat handlungsfähig und vertrauenswürdig ist, muss er stets die vollständige Hoheit über seine IT-Systeme und seine Daten behalten. Dafür sorgt Dataport mit hohen Sicherheitsstandards, klaren Datenschutzvorgaben und vorausschauender Unternehmenspolitik. Damit Bürger*innen und Unternehmen der Verwaltung ihre Daten anvertrauen können.
\vspace{1.0\baselineskip}

\noindent
{\em Kooperation:}\\
Seit der Gründung 2004 organisiert Dataport mit seinen Trägern erfolgreiche IT-Kooperationen. Der Dienstleister sorgt für Kooperationen auf verschiedenen Ebenen. Durch gemeinsam genutzte Infrastrukturen wie das Twin Data Center. Durch die gemeinsame Entwicklung von IT-Lösungen. Durch das Schaffen von Schnittstellen in IT-Systemen, die eine Zusammenarbeit ermöglichen. Länderübergreifend zwischen Bundesländern sowie ebenenübergreifend zwischen Bund, Ländern und Kommunen.

\newpage
\include{sponsorentexte/004-technologiestiftung_berlin}
\newpage
\section*{Goldponsor und Aussteller}
\begin{flushright}
\includegraphics[width=0.7\textwidth]{101_BKG_Logo_RGB.png}
\end{flushright}
\noindent
Das {\bfseries Bundesamt für Kartographie und Geodäsie (BKG)} ist eine Behörde im Geschäftsbereich des Bundesministeriums des Innern und für Heimat (BMI). Es fungiert als zentraler Dienstleister des Bundes und Kompetenzzentrum für Geoinformation und geodätische Referenzsysteme. Das BKG befasst sich mit der Beobachtung sowie der Datenhaltung bis hin zur Analyse, Kombination und Bereitstellung von Geodaten. Das BKG ermöglicht aufgrund der zentralen Geodatenbereitstellung eine optimale und wirtschaftliche Geodatennutzung im Bundesbereich.

\noindent
Das BKG setzt sich für eine offene Datenpolitik ein, wodurch die Verbreitung von Open Data gefördert wird. Dies schließt die Beratung anderer Bundesbehörden beim Umgang mit OSM-Daten ein. Die Nutzung, Entwicklung und Verbreitung der Nutzung freier Software liegen ebenfalls im Bereich der Aktivitäten.

\noindent
Von der Arbeit des BKG profitieren insbesondere Bundeseinrichtungen, die öffentliche Verwaltung, Wirtschaft, Wissenschaft~-- und fast jeder Bürger in Deutschland. Experten aus den verschiedensten Bereichen wie Verkehr, Katastrophenvorsorge, Innere Sicherheit, Energie und Umwelt verwenden Geodaten, Landkarten, Referenzsysteme und Informationsdienste des BKG für ihre Pläne und Untersuchungen. Das BKG unterhält ein Dienstleistungszentrum in Leipzig sowie geodätische Observatorien im In- und Ausland.

\newpage
\include{sponsorentexte/102-opengis}
\normalsize

